\subsection*{A source-gated final A.11 ratio and an integration-data obstruction}
\label{subsec:a11-source-gated-ratio-20261021}
On both A.11 correction supports the source has $U=M=J=0$.
Its exact stress identity (4.16) consequently reads
\[
N_s=\frac{4h\eta S-(1-\eta^2)S_\eta}{X}
+4A\eta\Pi-(1-\eta^2)\Pi_\eta,\qquad A=\tfrac12+h.
\]
The formulas
$S=\frac12\int_X^\infty E^2\,dx$ and
$\Pi=-\frac12\int_X^\infty E^2\,dx/x$
require, respectively, the true source conditions
$S(\infty)=\Pi(\infty)=0$. They are not implied by the
two A.11 moment-correction equations alone.
With $D_s=S/(XE^2)$, $P_s=-\Pi/E^2$, nonpositive
future logarithmic slopes and
$\sup_{x\ge X}|\partial_\eta\log E(x,\eta)|\le\varepsilon$,
the source identities yield
\[
0\le P_s\le\tfrac12,\qquad
|N_s|/E\le E\{4hD_s+4AP_s+2\varepsilon(D_s+P_s)\}.
\]
\begin{proposition}[Sharpened exact-rational sufficient stress gate]
\label{prop:a11-final-gate-20261021}
Suppose on BOTH complete A.11 supports and every $|\eta|\le1$
that the above genuine source tails hold and
\[
0<\lambda\le\tfrac1{20},\quad0<h\le\lambda/100,\quad
0<E\le\lambda/4,\quad0\le D_s\le233/100,\quad
0\le P_s\le1/2,\quad
\varepsilon\le\frac{2560\lambda^{28}}{1-40960\lambda^{28}}.
\]
Then $|N_s|/E<(503/2000)\lambda<\lambda/2$.
If also $Q_s\ge(49/50)\lambda$, $X\ge9/\lambda$,
$-3\lambda/2\le\ell\le-\lambda/2$ and
$0<L=1-2h\eta^2\le1$, then with
$a=2-2\ell$, $p_1=XQ_s/L$ and $w=N_s/(EQ_s)$,
\[
2(1-a/p_1)^2-(a-2)w^2>
\frac{85676973}{76832000}>\frac9{20}.
\]
\end{proposition}
\begin{proof}
The tail bracket is bounded by
\[
\frac{233}{50000}+1+\frac1{1000}
+\frac{283}{50}\frac{2560\,20^{-28}}{1-40960\,20^{-28}}
<\frac{503}{500}.
\]
Multiplication by $E\le\lambda/4$ proves the first claim.
The other hypotheses give
$|w|<503/1960$, $p_1\ge441/50$ and
$2\le a\le43/20$, so $a/p_1<1/4$.
The displayed cone bound follows from
$2(3/4)^2-(3/20)(503/1960)^2
=85676973/76832000$.
\end{proof}
\begin{proposition}[Why terminal normalization is indispensable]
\label{prop:a11-gauge-obstruction-20261021}
The A.11 bump equations alone do not imply any uniform
bound on $|N_s|/E$ unless the source's accumulated moment
and pressure data are fixed. Formally replacing
$S$ by $S+C_S(\eta)$ and $\Pi$ by $\Pi+C_\Pi(\eta)$
changes $N_s$ by
\[
\Delta N_s=
\frac{4h\eta C_S-(1-\eta^2)C_S'}{X}
+4A\eta C_\Pi-(1-\eta^2)C_\Pi'.
\]
For $C_S=K\eta$, $C_\Pi=0$, this becomes
$\Delta N_s(0)=-K/X$, unbounded as $|K|\to\infty$.
\end{proposition}
This algebraic separation is NOT a counterexample to the
complete source, whose moments and pressure have fixed physical
integral definitions. It proves why genuine terminal gauges and
one independently authenticated A.11 interval covering are
necessary proof inputs. The strict arithmetic gate alone
does not certify any actual source tuple, full exterior cone,
Appendix-B contraction, unforced startup system or
Theorem~1.1. The accompanying finite Lean statements likewise
do not formalize the source integral/PDE arguments.
